\chapter[1988 --- Black et al.]{1988: James W. Black, Gertrude B. Elion, George H. Hitchings}
\label{ch:1988_Black}

\section*{Main Contribution}

\textit{Discovered important principles for drug treatment---specifically, that drugs can be designed to block specific receptors---founding the field of rational drug design.}\cite{nobel-1988,lecture-1988}

\section{Official Citation}

for their discoveries of important principles for drug treatment

The 1988 Nobel Prize in Physiology or Medicine was awarded jointly to James W. Black, Gertrude B. Elion, and George H. Hitchings for their discoveries of important principles for drug treatment. Black's development of beta-adrenergic receptor blockers (beta-blockers) and histamine H$_2$-receptor antagonists for the treatment of peptic ulcer disease, and Elion and Hitchings' rational design of antimetabolite drugs for the treatment of cancer, gout, and infectious diseases, revolutionized the practice of pharmacology and established the principle that drugs can be designed to interact with specific molecular targets in the body.

\section{Background and Historical Context}

For most of the history of medicine, drug discovery was a haphazard process. New drugs were discovered by trial and error---often by accident---and their mechanisms of action were unknown. Aspirin, discovered in 1897, was used effectively for decades before its mechanism of action (inhibition of cyclooxygenase) was elucidated. Penicillin, discovered by Alexander Fleming in 1928, was used to treat bacterial infections for years before its mechanism of action (inhibition of cell wall synthesis) was understood. The idea that drugs could be designed rationally---based on a detailed understanding of the molecular targets they interact with---was a significant concept that emerged in the mid-twentieth century.

James Whyte Black was born in Uddingston, Scotland, in 1924, and trained in medicine at the University of St. Andrews. After a period of clinical practice, Black turned to pharmacological research, initially at the ICI (Imperial Chemical Industries) pharmaceutical division in England and later at the Wellcome Research Laboratories and the University of London.

Gertrude Belle Elion was born in New York City in 1918 and trained in chemistry at Hunter College. She was unable to pursue graduate studies because of the Great Depression and the gender discrimination prevalent at the time, so she worked as a laboratory assistant before joining the Wellcome Research Laboratories in Tuckahoe, New York, in 1944.

George Hitchings was born in Hoquiam, Washington, in 1905, and trained in chemistry and medicine at the University of Washington and Harvard University. He joined the Wellcome Research Laboratories in 1942 and became the director of the Department of Experimental Therapy, where he supervised Elion's work.

\section{The Discovery/Contribution}

\subsection{James W. Black and the Development of Beta-Blockers}

Black's first major contribution was the development of propranolol, the first clinically successful beta-adrenergic receptor blocker. In the late 1950s, Black recognized that the sympathetic nervous system---which mediates the body's ``fight or flight'' response---exerts its effects through specific receptors on the surface of target cells. These receptors, which he called beta-adrenergic receptors, are activated by the hormones epinephrine (adrenaline) and norepinephrine (noradrenaline) and mediate a range of physiological responses, including increased heart rate, increased cardiac contractility, and dilation of the airways.

Black hypothesized that a drug that could selectively block beta-adrenergic receptors would be useful for the treatment of angina pectoris (chest pain caused by reduced blood flow to the heart) and other cardiovascular diseases. Working with the ICI pharmaceutical company, Black developed propranolol, a beta-blocker that was approved for clinical use in 1964. Propranolol was the first drug to provide effective relief from angina pectoris, and it quickly became one of the most widely prescribed cardiovascular drugs in the world.

Propranolol works by blocking the beta-adrenergic receptors on the surface of heart muscle cells, thereby reducing the stimulatory effects of epinephrine and norepinephrine. This results in a decreased heart rate, decreased cardiac contractility, and reduced oxygen demand by the heart---all of which help to relieve angina and reduce the risk of heart attack. Propranolol also proved to be effective for the treatment of hypertension (high blood pressure), cardiac arrhythmias (abnormal heart rhythms), and migraine headaches.

The development of propranolol was based on the principle of ``rational drug design''---the idea that a drug could be designed to interact with a specific molecular target (in this case, the beta-adrenergic receptor) based on a detailed understanding of the target's structure and function. This principle was significant at the time, and it has since become the dominant paradigm in drug discovery.

Following the success of propranolol, many other beta-blockers were developed, including atenolol, metoprolol, bisoprolol, and carvedilol. These drugs differ in their pharmacological properties (such as selectivity for beta-1 vs. beta-2 receptors, and the presence of additional pharmacological activities) but share the basic mechanism of blocking beta-adrenergic receptors. Beta-blockers are now among the most widely used drugs in the world and are a cornerstone of the treatment of cardiovascular disease.

\subsection{James W. Black and the Development of H\texorpdfstring{$_2$}{2}-Receptor Antagonists}

Black's second major contribution was the development of cimetidine (Tagamet), the first histamine H$_2$-receptor antagonist, for the treatment of peptic ulcer disease. In the 1960s, Black recognized that histamine---a signaling molecule produced by the body---acted through at least two distinct types of receptors. The H$_1$-receptor mediates the allergic responses to histamine (such as hay fever and asthma), and these responses could be blocked by antihistamines such as diphenhydramine (Benadryl). However, the H$_2$-receptor, which mediates the stimulation of gastric acid secretion by histamine, could not be blocked by conventional antihistamines.

Black hypothesized that a drug that selectively blocked H$_2$-receptors would inhibit gastric acid secretion and thereby provide an effective treatment for peptic ulcer disease---a common and often debilitating condition that affected millions of people worldwide. Working with a team of chemists and pharmacologists at the Smith Kline and French Laboratories (now GlaxoSmithKline), Black developed cimetidine, which was approved for clinical use in 1976.

Cimetidine works by blocking the H$_2$-receptors on the surface of parietal cells in the stomach, thereby reducing the production of gastric acid. This promotes the healing of peptic ulcers and provides relief from the symptoms of gastroesophageal reflux disease (GERD). The development of cimetidine was a landmark in the treatment of peptic ulcer disease, which had previously been managed primarily by dietary restrictions and, in severe cases, by surgery (vagotomy and partial gastrectomy). Cimetidine and its successors (ranitidine/famotidine) transformed the management of peptic ulcer disease from a surgical to a medical condition, and they remain among the most widely used gastrointestinal drugs in the world.

\subsection{Gertrude Elion and George Hitchings and the Rational Design of Antimetabolite Drugs}

Elion and Hitchings' contribution was the development of a rational approach to drug design based on the differences in nucleic acid metabolism between normal cells and pathogenic organisms. Their approach was based on the principle that cells which divide rapidly---such as cancer cells, bacteria, and protozoa---require large quantities of nucleotides (the building blocks of DNA and RNA) and are therefore more susceptible to drugs that interfere with nucleotide synthesis than are slowly dividing normal cells.

Working at the Wellcome Research Laboratories, Elion and Hitchings developed a series of antimetabolite drugs---molecules that mimic the structure of natural nucleotides and interfere with the enzymes involved in nucleotide synthesis. Their first major success was 6-mercaptopurine (6-MP), a purine analog that inhibits the synthesis of nucleic acids in rapidly dividing cells. 6-MP was one of the first effective drugs for the treatment of acute leukemia in children, and it dramatically improved the survival of children with this previously fatal disease.

Elion and Hitchings went on to develop a series of other antimetabolite drugs, including:

\begin{itemize}
 \item \textbf{Azathioprine}, a prodrug of 6-mercaptopurine that is used as an immunosuppressive agent in organ transplantation and autoimmune diseases. Azathioprine was one of the first drugs to make organ transplantation feasible by suppressing the immune rejection of transplanted organs.
 \item \textbf{Allopurinol}, a xanthine oxidase inhibitor that is used to treat gout and hyperuricemia. Allopurinol works by inhibiting the enzyme that converts purines to uric acid, thereby reducing the levels of uric acid in the blood and preventing the formation of urate crystals in the joints.
 \item \textbf{Pyrimethamine}, an antifolate drug that inhibits the enzyme dihydrofolate reductase in protozoa. Pyrimethamine is used to treat and prevent malaria and toxoplasmosis, and it works by blocking the synthesis of tetrahydrofolate, a cofactor required for the synthesis of nucleic acids.
 \item \textbf{Trimethoprim}, another dihydrofolate reductase inhibitor that is used to treat bacterial urinary tract infections and other bacterial infections. Trimethoprim is often used in combination with sulfamethoxazole (as co-trimoxazole/Bactrim), and the combination is highly effective against a wide range of bacterial pathogens.
\end{itemize}

The approach that Elion and Hitchings developed---the rational design of drugs based on the differences in metabolism between target organisms and host cells---was a pioneering example of what is now called ``targeted therapy.'' Their work demonstrated that drugs could be designed to exploit specific biochemical differences between pathogenic and normal cells, thereby achieving selective toxicity with minimal side effects.

\section{Impact on Modern Medicine}

The impact of Black, Elion, and Hitchings on modern medicine has been enormous and multifaceted. Black's development of beta-blockers and H$_2$-receptor antagonists has improved the lives of hundreds of millions of people with cardiovascular disease and peptic ulcer disease. Beta-blockers are now among the most widely prescribed drugs in the world, with an estimated 100 million prescriptions written annually in the United States alone. They are used to treat hypertension, angina, heart failure, cardiac arrhythmias, and migraine headaches, and they have been shown to reduce the risk of heart attacks and cardiovascular death.

The development of H$_2$-receptor antagonists (and, later, proton pump inhibitors) transformed the management of peptic ulcer disease from a surgical to a medical condition. Before the introduction of cimetidine, peptic ulcer disease was a leading cause of hospitalization and surgery, and it affected millions of people worldwide. Today, peptic ulcer disease is effectively managed with medications, and the need for surgery has been dramatically reduced.

Elion and Hitchings' antimetabolite drugs have saved millions of lives. 6-Mercaptopurine and its derivatives remain important components of chemotherapy regimens for acute leukemia and other cancers. Azathioprine is a cornerstone of immunosuppressive therapy in organ transplantation and autoimmune diseases. Allopurinol is the most widely used drug for the treatment of gout. Trimethoprim-sulfamethoxazole is one of the most commonly prescribed antibiotics for the treatment of urinary tract infections and other bacterial infections.

More broadly, the work of the three laureates established the principle that drugs can be designed rationally, based on a detailed understanding of the molecular targets they interact with. This principle has become a major paradigm in drug discovery, and it has led to the development of many important drugs of the past half-century, including statins, ACE inhibitors, angiotensin receptor blockers, and targeted cancer therapies.

\section{Legacy}

James Black continued his work on drug design at the Wellcome Foundation and later at the University of London, where he became a professor and director of the James Black Foundation. He was knighted by Queen Elizabeth II in 1981 and received the Order of Merit in 2000. Black died in 2010, but his contributions to pharmacology continue to influence drug discovery and development.

Gertrude Elion continued her work at the Wellcome Research Laboratories and later at the Duke University Medical Center, where she served as a research professor and adviser. She was the first woman to be inducted into the National Inventors Hall of Fame (in 1991) and received the National Medal of Science (in 1993). Elion died in 1999, but her antimetabolite drugs continue to save lives around the world.

George Hitchings continued to lead research at the Wellcome Laboratories until his retirement in 1975. He died in 1998, but his legacy lives on in the drugs he helped to develop and in the rational approach to drug design that he and Elion pioneered.

The work of the three laureates of 1988 demonstrated that drug discovery need not be a matter of chance or trial and error. By understanding the molecular targets of drugs and the biochemical differences between normal and pathological states, it is possible to design drugs that are more effective, more selective, and safer than those discovered by traditional methods. This principle has transformed the pharmaceutical industry and has led to the development of some of a major and life-saving drugs in the history of medicine.


\section{Modern Developments}

Black, Elion, and Hitchings' rational drug design work has led to modern computational drug discovery. Structure-based drug design uses X-ray crystallography and cryo-EM to design drugs targeting specific protein structures. AI-driven drug discovery (AlphaFold for protein structure prediction) accelerates the drug development pipeline. Understanding of receptor pharmacology has led to targeted therapies for cancer, cardiovascular disease, and neurological disorders.


\section{Bibliography}

\begin{thebibliography}{9}
\bibitem{nobel-1988}
Nobel Prize Outreach, ``The Nobel Prize in Physiology or Medicine 1988,''\emph{NobelPrize.org}.\\
\url{https://www.nobelprize.org/prizes/medicine/1988/summary/}

\bibitem{lecture-1988}
James W. Black, ``Nobel Lecture,''\emph{NobelPrize.org}. Winner-authored primary source; downloadable from the official Nobel Prize archive.\\
\url{https://www.nobelprize.org/prizes/medicine/1988/black/lecture/}
\end{thebibliography}
